\title{LongRoPE: Extending LLM Context Window Beyond 2 Million Tokens}

\begin{abstract}
A large context window is a sought-after characteristic in large language models (LLMs). Nevertheless, challenges such as the high expense of fine-tuning, limited availability of lengthy textual data, and the emergence of problematic values from novel token positions restrict today's extended context windows to a maximum of approximately 128k tokens.
\end{abstract}

In this work, we present LongRoPE, which is the first method to expand the context window of pre-trained LLMs to an impressive \textbf{2048k} tokens. This is accomplished with as few as 1k fine-tuning steps using training sequences up to 256k tokens, without compromising performance on the original short context window. LongRoPE builds on three principal contributions: (i) we discover and leverage two distinct non-uniformities in positional interpolation, utilizing an efficient search to provide improved initialization for fine-tuning, and allowing an 8$\times$ context length extension even without further fine-tuning; (ii) we propose a staged expansion approach, where the model is first fine-tuned at 256k token length and subsequently undergoes a second positional interpolation, applied to the newly extended model, to reach the full 2048k context window; (iii) we further recalibrate LongRoPE on sequences of 8k tokens to restore high accuracy at shorter context windows. Thorough evaluations on LLaMA2 and Mistral models across a range of tasks confirm the superior performance of our approach. LongRoPE-extended models preserve the original neural network design with only minor changes to the positional encoding mechanism, enabling compatibility with most existing optimizations. The implementation will be released at \texttt{https://github.com/microsoft/LongRoPE}
\end{abstract}

\section{Introduction}

Large Language Models (LLMs) have achieved significant milestones across numerous applications~\cite{gpt4,llama2}, yet they are constrained by relatively short context windows, such as the 4096-token boundary in LLaMA2~\cite{llama2}. When processing input that exceeds this window, model effectiveness diminishes because it encounters positional inputs absent from its training distribution. This constraint becomes particularly problematic in key use cases such as in-context learning with large sets of examples~\cite{Huang2023FewerIM}, as well as in the deployment of LLM agents~\cite{park2023generative,madaan2023selfrefine}.

Recent studies have demonstrated that it is possible to expand a pre-trained LLM's context window up to approximately 128k tokens through fine-tuning with longer sequences~\cite{longlora,pi,yarn,zhang2024soaring,liu2023scaling}. However, scaling the context window beyond this size presents several significant challenges. The first issue arises from introducing previously unseen positional indices during extension, which can generate numerous outlier values and result in out-of-distribution complications, thereby hindering fine-tuning convergence~\cite{pi}. This difficulty becomes even more pronounced when enlarging the window from 4k to over 1000k, as this process can involve introducing more than 90\% new positions. A second challenge relates to the requirement of matching long text samples for effective fine-tuning. Currently available datasets rarely contain a sufficient number of texts longer than 1000k tokens. In addition, training on such lengthy sequences is computationally demanding, necessitating substantial GPU resources and an impractical amount of training time. Finally, as the context window grows to extreme lengths, attention mechanisms are forced to distribute focus across a vastly greater number of token positions, leading to attention dilution and deteriorating the model’s effectiveness on shorter, original contexts~\cite{pi}.

A common strategy for addressing the initial challenge involves adjusting RoPE positional embeddings~\cite{rope,pi} by mapping extended position indices back into the pretrained positional range, as illustrated in Fig.~\ref{fig:overview}. The Position Interpolation (PI) technique~\cite{pi} achieves this by performing a linear interpolation of RoPE’s rotary angles in proportion to the desired extension. Meanwhile, NTK~\cite{ntk,dynamicntk} proposes a non-uniform scheme that applies different interpolation and extrapolation strategies to various RoPE dimensions. YaRN~\cite{yarn}, in contrast, sorts RoPE dimensions into three groups based on frequency and applies extrapolation, NTK, or linear interpolation accordingly. Nevertheless, the nature of positional embeddings within Transformer architectures is characterized by \emph{complex, non-uniform information entropy}. Existing interpolation methods typically overlook this intricate non-uniformity, which can cause a loss of information and ultimately restricts the maximum attainable context window.

Section~\ref{sec:analysis} presents empirical evidence supporting two main conclusions: \textbf{(1)} Successful positional interpolation requires addressing two distinct non-uniformities—differences among RoPE dimensions and disparities in token positions. Less interpolation is needed for lower RoPE dimensions and earlier token positions, although the most effective approach varies depending on the length one aims to extend to. \textbf{(2)} By explicitly incorporating these non-uniformities into the positional interpolation process, it becomes possible to better preserve crucial information from the original RoPE, especially regarding important dimensions and token locations. This approach significantly reduces information loss incurred during interpolation and leads to improved initial states for subsequent fine-tuning. Furthermore, in scenarios where models are not fine-tuned, this method enables sequence length extensions up to 8$\times$ the original without substantial degradation.

Building upon these results, we propose \textbf{LongRoPE}, a powerful technique capable of expanding the context window of LLMs to more than 2 \emph{million} tokens. The design of LongRoPE incorporates three major advancements. The first of these leverages multidimensional non-uniformities within positional interpolation to their fullest potential. By determining optimal rescale factors for RoPE's rotation angles specific to each RoPE dimension, and informed by the respective token positions, LongRoPE enhances performance. Given that the parameter space for rescale factor selection grows exponentially relative to the desired extension ratio, LongRoPE employs an evolutionary search algorithm coupled with two optimization strategies to significantly improve search efficiency. An illustration of the resultant rescaled RoPE configuration derived from this search is provided in Fig.~\ref{fig:overview}.

Subsequently, {LongRoPE} adopts a progressive and efficient extension method to reach a 2048k context window, circumventing the necessity to fine-tune directly on exceedingly lengthy texts—which are both rare and difficult to obtain. The process initiates with identifying an optimal 256k length using the pre-trained LLM, followed by fine-tuning at this window size. Since the non-uniform positional interpolation technique supports an 8$\times$ expansion in settings without additional fine-tuning, a secondary search for appropriate RoPE rescale factors is performed on the previously fine-tuned, extended LLM. This process results in achieving a 2048k context window for LLaMA2 and Mistral~\cite{mistral}.

In order to address the potential drop in performance on the original, shorter context windows, {LongRoPE} further fine-tunes the RoPE rescale factors on the expanded LLM. Analogous to the scaling up process from 256k to 2048k, we utilize our search algorithm to rescale down to 4k and 8k context lengths on the 256k fine-tuned LLM, which reduces reliance on positional interpolation. At inference time, for sequences shorter than 8k tokens, we apply the optimized rescale factors to RoPE.

Comprehensive evaluations performed on multiple LLMs and a range of long-context tasks validate the efficacy of our approach. Our results indicate that LongRoPE consistently preserves low perplexity across evaluation lengths spanning from 4k to 2048k, attains passkey retrieval accuracy exceeding 90\%, and achieves accuracy comparable to baseline models on standard benchmarks situated within the 4096 context window. Notably, LongRoPE is broadly applicable to any LLM utilizing RoPE embeddings. We plan to make our code and the LongRoPE-2048k models publicly available.

\section{Non-uniformity in Positional Interpolation}
\label{sec:analysis}

\subsection{Preliminary}

Transformer architectures depend on explicit positional encodings, commonly implemented as position embeddings, to capture sequence order among input tokens. In this study, we concentrate on RoPE~\cite{rope} positional embeddings, a method extensively adopted in contemporary LLMs. For a token located at position $n$, the associated RoPE embedding can be described more succinctly as:

\begin{equation}
		\fontsize{7.8}{7.8} \selectfont
	\label{eq:rope}
	\left[ cos(n\theta_0), sin(n\theta_0), cos(n\theta_1), \cdot\cdot\cdot, cos(n\theta_{d/2-1}), sin(n\theta_{d/2-1}) \right]   
\end{equation}

In this formulation, $d$ stands for the embedding dimension. The rotary position angle for a token at position $n$ is given by $n\theta_i$, where each $\theta_i$ determines a rotation frequency as $\theta_i = \theta^{-2i/d}$. Within RoPE, the typical base value assigned to $\theta$ is 10000.

\textbf{\emph{Context window extension ratio $s$} and positional interpolation}. We define $s$ as  the ratio of extended context length $L'$ to the original length $L$: $s=\frac{L'}{L}$. 

In efforts to increase the context window from $L$ to $L'$, prevailing positional interpolation techniques recommend decreasing rotation frequencies $\theta_i$ proportionally to the extension ratio $s$. Define $\beta=\theta^{2/d}$ and let $\lambda$ be the effective rescaling factor as determined by $s$. These approaches to positional interpolation can thus be represented in a unified framework as:

\begin{equation}	
	\fontsize{7.5}{7.5} \selectfont
	\label{eq:unified_rope}
	\begin{aligned}
		\left[cos\left(\frac{n}{\lambda(\beta)^0}\right), sin \left(\frac{n}{\lambda(\beta)^0}\right), cos\left(\frac{n}{\lambda(\beta)^1}\right), \cdot\cdot\cdot,  sin \left(\frac{n}{\lambda(\beta)^{d/2-1}}\right) \right]   
	\end{aligned}
\end{equation}

\textbf{Linear positional interpolation (PI)}. The PI method~\cite{pi} employs linear scaling of position indices confined to the original pre-training sequence length. Specifically, when extending to a new sequence length defined by a scale factor $s$, the model uniformly compresses all rotation angles by a factor of $\lambda=s$ along each dimension of RoPE. This uniform compression leads to overly compacted positional representations, which diminishes the model's capacity to discern tokens with nearby positions. Consequently, PI often exhibits degraded performance when applied at larger extension ratios.

\textbf{NTK-based interpolation and extrapolation}. ~\cite{ntk,dynamicntk} examine RoPE through the framework of information encoding and employ Neural Tangent Kernel (NTK) theory~\cite{jacot2018neural,tancik2020fourier} in their analysis. To address the issue of position crowding in PI, they propose redistributing the interpolation burden among different RoPE dimensions. Specifically, the approach applies weaker scaling to dimensions associated with higher frequencies and stronger scaling to those with lower frequencies. This design enables both positional interpolation and extrapolation by setting $\lambda=s^{i}$. The enhanced dynamic NTK method~\cite{dynamicntk} further refines the extension ratio at each position in accordance with the actual sequence length. In contrast to PI, which relies on fine-tuning, NTK-aware approaches can expand the context window without requiring fine-tuning. However, such extensions are commonly capped at a maximum ratio of 4$\times$.

\textbf{YaRN}~\cite{yarn} presents a notable advancement in positional interpolation by categorizing RoPE dimensions into three frequency-specific segments, each assigned its own interpolation technique. Specifically, extrapolation ($\lambda$=1) is applied to high-frequency dimensions, while low-frequency dimensions are interpolated linearly (PI). Those dimensions at intermediate frequencies are addressed using the NTK approach. The effectiveness of YaRN stems from its partitioning of RoPE dimensions; however, this grouping is currently determined by empirical human judgment, potentially leading to less-than-optimal outcomes when adapting to novel LLMs.

\subsection{Study on Non-uniform Positional Interpolation}

Drawing on insights from NTK and YaRN, we observe that their improvements are largely attributed to leveraging non-linearities, particularly by utilizing varying frequency treatments across RoPE dimensions for advanced interpolation and extrapolation. Nonetheless, the existing approaches to non-linearity predominantly depend on hand-crafted human heuristics. This situation prompts two central questions: (1) Is the present approach to positional interpolation the best achievable? (2) Might there exist other, yet unidentified, forms of non-linearity that could be beneficial?

To address these questions, we employ evolution search (refer to Sec.~\ref{sec:method}) to identify improved non-uniform positional interpolation schemes for LLaMA2-7B. The evolution process is steered by evaluating perplexity, utilizing 5 randomly selected examples from the PG19~\cite{pg19} validation dataset. Our experimental investigation leads to several significant observations.

\textbf{Finding 1}: \textit{RoPE dimensions exhibit substantial non-uniformities, which are not effectively handled by current positional interpolation methods.}

We conduct a search for the optimal value of $\lambda$ for each RoPE dimension as outlined in Eq.~\ref{eq:unified_rope}. Table~\ref{tbl:exp1} presents the perplexity outcomes for LLaMA2-7B across various approaches on the PG19 and Proof-pile~\cite{proof-pile} test sets, evaluated without any fine-tuning. The approach derived from our search delivers clear improvements, indicating that both the linear (PI) and the non-uniform interpolation methods (Dynamic-NTK and YaRN) are not fully optimal. Importantly, YaRN performs worse than both PI and NTK on PG19, as it fails to extend to the desired context window length when the LLM has not been fine-tuned. Specifically, for an 8k context window, the perplexity with YaRN increases sharply beyond 7k.

In our investigation, we found that the rescaling factors $\lambda$ in Eq.~\ref{eq:unified_rope} vary across dimensions, unlike the uniform scaling factor $s$ employed in PI, the calculation method in NTK, or the group-based computation in YaRN. The introduction of these dimension-specific scaling factors leads to a notable enhancement in LLaMA2’s language modeling capability—as measured by perplexity—over extended context windows of 8k and 16k, all without any fine-tuning. This improvement stems from the modified positional embedding’s ability to largely maintain the structure of the original RoPE, particularly in key dimensions, thereby alleviating the model’s challenge in discriminating among closely positioned tokens.

\textbf{Finding 2}: \textit{RoPE for the initial tokens in the input sequence should be extrapolated with less interpolation.} 

We propose that for the initial $\hat{n}$ tokens in the input sequence, RoPE should incorporate less interpolation. The rationale behind this is that these tokens tend to receive higher attention scores, rendering them especially important in the attention mechanism, as noted in Streaming LLM~\cite{streamingllm} and LM-Infinite~\cite{han2023lminfinite}. To test this hypothesis, we expanded the context window to 8k and 16k tokens utilizing PI and NTK, while selectively applying no interpolation to the first $\hat n$ tokens (where $\hat n$ = 0, 2, ..., 256). When $\hat n=0$, the setup corresponds to the standard PI and NTK configuration. Table~\ref{tbl:tokenposition} offers two main findings: \textbf{(1)} omitting position interpolation for the initial tokens leads to enhanced performance; and \textbf{(2)} the most effective value for $\hat n$ varies according to the length of the context window being extended.

\textbf{Finding 3}: \textit{Non-uniform positional interpolation effectively extends LLM context window in both fine-tuning and non-fine-tuning settings.} 

Although our results demonstrate that the optimized non-uniform position interpolation notably enhances extension capabilities at 8k and 16k context lengths without the need for fine-tuning, extending to even longer sequences necessitates model adaptation. Accordingly, we adapt LLaMA2-7B utilizing the discovered RoPE configuration, specifically targeting a 64k context window (refer to the Appendix for detailed experimental settings). As presented in Table~\ref{tbl:finetune}, our approach delivers substantial improvements compared to PI and YaRN, both in the original and fine-tuned LLaMA2-7B variants. This advantage stems from our precise exploitation of non-uniform positional interpolation, which reduces information degradation and yields a superior initialization prior to fine-tuning.

\textbf{Summary}. This work reveals two distinct types of non-uniformity—those associated with varying RoPE dimensions and with token positions. Leveraging these forms of non-uniformity within positional interpolation leads to marked advancements in the extension of LLM context windows.

\section{LongRoPE}

\label{sec:method}
Building on these observations, we propose LongRoPE, an approach that initially deploys an effective search strategy to comprehensively leverage both types of non-uniformities. Utilizing this algorithm, we are able to expand the LLM context window to exceed 2 million tokens.

\subsection{Problem Formulation}

The presence of these two non-uniformities results in an extensive solution space and increases the complexity of the optimization process. To tackle this challenge, we recast the multidimensional non-uniform position interpolation optimization as a search problem.

Consider an LLM designed for a context window of size $L'$, processing long input sequences $\mathbf{X}$, each $\mathbf{x} \in \mathbf{X}$ exceeding $L'$ tokens. Let the canonical rotary angle for the $i^{th}$ dimension in the RoPE embedding at token index $n$ be represented as $\frac{n}{\beta^i}$. The corresponding optimization task can therefore be formalized as:

\begin{equation}
\begin{gathered}
		\label{eq:problem}

\underset{ \mathbf{x} \in \mathbf{X}; \, |\mathbf{x}| \geq L' }{ \operatorname{arg\,min} } \; \mathcal{L}\, \left( \text{LLM} (\text{RoPE}, \mathbf{X}) \right), \quad \text{where} \;
\\
\scalebox{0.93}{
$\underset{ \begin{subarray}{l} i = 0, \ldots, \frac{d}{2} - 1; \\ n \in [0, |\mathbf{x}|); \end{subarray} }{ \text{RoPE}_{(n)} } = \left[ \ldots, \cos \left( \mathbb{I}(\hat{\lambda}_{i}, \hat{n}) \cdot \frac{n}{\beta^i} \right), \sin \left( \mathbb{I}(\hat{\lambda}_{i}, \hat{n}) \cdot \frac{n}{\beta^i} \right), \ldots \right]$}
\\
\text{where} \; \mathbb{I}(\hat{\lambda}_{i},\hat{n}) = 
\begin{cases}
1 & \text{if } n < \hat{n} \\
\frac{1}{\lambda_i} & \text{if } n \geq \hat{n}
\end{cases}
\end{gathered}
\end{equation}

In this context, we define a collection of rescale factors, $\mathbb{I}(\hat{\lambda}_{i},\hat{n})$, to address the two types of non-uniformities present. Here, $\hat{\lambda_i}$ and $\hat{n}$ signify the non-uniform characteristics related to RoPE dimensions and to the positions of tokens, correspondingly. More precisely, the rotation angle in the $i^{th}$ RoPE dimension is adjusted by $\mathbb{I}(\hat{\lambda}_{i},\hat{n})$, with $\hat\lambda_{i}$ serving as the adjustment parameter and $\hat{n}$ as the token position cut-off. When dealing with the first $\hat{n}$–1 token positions, the rescale factor $\hat\lambda_{i}$ is not yet implemented, thus the standard RoPE angle $\frac{n}{\beta^i}$ remains unchanged. The adjustment through the rescale factor becomes effective for all token positions where $n\ge\hat{n}$.

Suppose we are given a desired context window of size $L'$. The goal is to determine the optimal set of rescale factors, denoted as $\mathbb{I}(\hat{\lambda}_{0},\hat{n}), \mathbb{I}(\hat{\lambda}_{1},\hat{n}), \ldots, \mathbb{I}(\hat{\lambda}_{i},\hat{n}), \ldots$, applied across the $1^{st}$ through $d^{th}$ RoPE dimensions. The intention is for the target $\text{LLM}$, equipped with these adjusted $\text{RoPE}$ parameters, to minimize the next token prediction loss $\mathcal{L}$ (that is, the perplexity) over input sequences $\mathbf{X}$ consisting of $L'$ tokens.

\subsection{Searching the Non-uniform Position Interpolation}

To address the issue outlined in Eq.~\ref{eq:problem}, we present a straightforward yet powerful approach that seeks optimal RoPE rescale factors. This strategy is designed to maximize the utilization of positional embedding’s multidimensional variations.

\textbf{Search space}. Our approach constructs an extensive search space that encapsulates both types of non-uniformities, as outlined in Table~\ref{tbl:searchspace}. Concretely, we permit the assignment of an individually optimized rescale factor to each dimension within the RoPE embedding. To facilitate the structure of the search space, our optimization focuses on $\lambda_i$ and $\hat{n}$, replacing direct exploration of $\mathbb{I}(\hat{\lambda}_{i},\hat{n})$, where  $\hat{\lambda}_{i}=1/\lambda_i$. As described in Table~\ref{tbl:searchspace}, each $\lambda_i$ is searched over a range that begins at a lower bound of 1.0 (corresponding to direct extrapolation) and extends up to $s\times1.25$ (permitting interpolation exceeding that of PI), with increments of 0.01. Here, $s$ denotes the ratio for extending the target context window.

The parameter $\hat{n}$ specifies how many of the initial token positions will use the unmodified RoPE embedding instead of undergoing position interpolation. In practice, $\hat{n}$ is selected from the set \{0, 1, 2, 4, 8, 12, 16, 20, 24, 28, 32, 64, 128, 256\}. If $\hat{n}=0$, then every token position utilizes the optimized rescale factors.

\textbf{Evolution-based search}. The search space outlined in Table~\ref{tbl:searchspace} contains a vast array of possible positional interpolation configurations, making thorough search extremely challenging. For instance, an extension factor of $s=4\times$ results in $400^{128/2}\times14=4\times10^{167}$ possible combinations. As the extension ratio increases, the search space grows at an exponential rate. To efficiently navigate this immense space, we employ evolution search~\cite{guo2020single} and incorporate two optimization strategies that considerably enhance search efficiency. The overall search workflow is depicted in Algorithm~\ref{alg:search}.

\textit{Optimized initial population generation}. In place of selecting all $P$ rescale factors purely at random, the initial population explicitly includes the three RoPE rescale factors associated with PI, NTK, and YaRN. The rest of the population, consisting of $P-3$ individuals, is then produced by applying random mutations—with probability $p$—to these three baseline rescale factors.

\textit{Monotonically non-decreasing constraint}. Once the initial population is created, each individual's perplexity is measured using the LLM. To do this, we implement the relevant RoPE rescale factors within the target LLM and then assess the perplexity on input $\mathbf{X}$. Those individuals with the top-k perplexity scores are chosen to serve as parents in the evolutionary process. Nevertheless, due to the expansive search space, straightforward crossover and mutation strategies often examine suboptimal candidates, resulting in a surplus of perplexity evaluations. This inefficiency becomes more pronounced as $L'$ increases, since each perplexity evaluation requires costly inference computations.

To mitigate this issue, we enforce a non-decreasing monotonic constraint on the sampled RoPE rescaling parameters: $\lambda_i\le \lambda_{i+1}$. Only those RoPE configurations meeting this requirement are utilized for evaluating perplexity in LLMs, which leads to a substantial reduction in search complexity. Concretely, we stipulate that $\lambda_i$ grows monotonically along with the RoPE dimension index ($i=0,\dots,63$). This form of dimensional monotonicity draws support from NTK theory~\cite{jacot2018neural,tancik2020fourier,ntk}, which indicates that RoPE’s lower dimensions—corresponding to higher frequencies—should employ less interpolation (i.e., have smaller $\lambda_i$ values), while higher dimensions—associated with lower frequencies—should employ greater interpolation (i.e., larger $\lambda_i$ values).

\begin{algorithm}[t]
	\small
	\caption{The search algorithm}
	\textbf{Input:} target LLM, input samples $\mathbf{X}$, population size $P$, mutation size $N_1$, crossover size $N_2$, maximum iterations $\mathcal{T}$, mutation probability $p$ \\

	\begin{algorithmic}[1]
		\label{alg:search}
		\STATE Initialize Top-k as $\phi$;
		\STATE $\text{P}_0$ $\gets$ \textit{Initialize\_population\_with\_optimization}($P$, $\mathbf{X}$, $p$); 
		\FOR{$i = 1$ {\bf to} $\mathcal{T}$}
		\STATE \textit{Compute\_perplexity}(LLM, $\text{P}_{i-1}$, $\mathbf{X}$);
		\STATE Top-k $\gets$ \textit{Update\_Topk}(Top-k, $\text{P}_{i-1}$);
		\STATE $\text{P}_{mutation}$ $\gets$ \textit{Mutation\_with\_mono\_constraint}(Top-k, $N_1$, $p$); 
		\STATE $\text{P}_{crossover}$ $\gets$ \textit{Crossover\_with\_mono\_constraint}(Top-k, $N_2$);
		\STATE $\text{P}_i$ $\gets$ $\text{P}_{mutation}$ $\cup$ $\text{P}_{crossover}$ $\cup$ Top-k;
		\ENDFOR
		\STATE Output the member in Top-k achieving the lowest perplexity;
	\end{algorithmic}
\end{algorithm}

\textbf{8$\times$ extension without fine-tuning}. Through evolutionary search, our approach efficiently discovers optimal non-uniform RoPE rescale factors, selectively maintaining crucial dimensions and positions to reduce the information loss that results from interpolation. As illustrated in Fig.\ref{fig:non-ft}, this technique allows the context window of LLaMA2 to be expanded from 4k to 32k tokens without the need for fine-tuning. In comparison, previous methods like PI, as well as non-uniform NTK and YaRN, exhibit a sharp increase in perplexity when extended beyond 2$\times$.

\subsection{Extending LLM Context Window to 2048K}
\label{sec:extendto2048k}

\textbf{Progressive extension to 2048k}. Here, we present our approach for expanding the context window of pre-trained LLMs from the standard 4k up to more than 2048k tokens. Our proposed non-uniform positional interpolation is capable of achieving an 8$\times$ expansion without the need for fine-tuning. However, when even greater extension is desired (such as 512$\times$), fine-tuning becomes necessary. One possible strategy involves determining appropriate RoPE rescaling factors at the intended 2048k context length and subsequently performing fine-tuning. Nonetheless, this approach is limited by the extremely high computational costs associated with such large-scale training. Additionally, in our experience, effectively fine-tuning LLMs under such substantial extension ratios presents significant difficulties (see Appendix).

Fortunately, LongRoPE demonstrates efficacy with both the base model and the fine-tuned variant at extended contexts. As a result, we propose a progressive, resource-efficient approach that allows us to reach the desired 2048k target using only 1k fine-tuning steps, all conducted within a training window of 256k length.

$\Diamond$ \textit{Scaling pre-trained LLMs to 256k via LongRoPE search}. Using LLaMA2 as a case study, we perform a search to accommodate context window sizes of 128k and 256k tokens. At these stages, the context window is expanded by factors of 32$\times$ and 64$\times$, respectively.

$\Diamond$ \textit{Fine-tuning to 256k}. Subsequently, we adapt the pre-trained LLM to support a context window of 256k tokens through further fine-tuning. Initially, LLaMA2 undergoes 400 fine-tuning steps utilizing RoPE rescaled factors corresponding to 128k. After this phase, we swap the RoPE rescaled factors to those configured for 256k, and continue fine-tuning for an additional 600 steps on the updated checkpoint. This staged approach demonstrates greater efficiency compared to commencing fine-tuning at 256k directly.

$\Diamond$ \textit{Scaling fine-tuned extended LLM to 2048k using LongRoPE search}. In the final step, we conduct an additional search process on the LLM that has already been fine-tuned for sequences of length 256k. Through this procedure, we achieve a remarkably large context window of 2048k tokens, all without the need for additional fine-tuning. This represents an extension by a factor of $512\times$.

\textbf{Recovery for shorter context windows}. Upon expanding to an extensive 2048k context window, we observe a decrease in model performance within the initial context range. This phenomenon is recognized in prior work on positional interpolation~\cite{pi}, where enforcing position embeddings into higher dimensions over the original window compresses these embeddings into a limited area, thereby harming the language model's capability. When the extension ratio reaches 512$\times$, the embeddings for positions in the initial 4k window are especially congested.

To address this issue, we conduct an additional evolutionary search over the expanded LLM, specifically tuning the RoPE rescale factors for shorter context lengths such as 4k and 8k. For these shorter sequences, we limit the upper bound for the searched $\lambda$ parameter, reflecting the decreased necessity for positional interpolation. At inference time, the LLM automatically modifies the relevant RoPE rescale factors in accordance with the context length.

\section{LongRoPE}

\label{sec:method}
Building on these insights, we propose LongRoPE, which incorporates a novel and efficient search algorithm designed to capitalize on both types of non-uniformities. Leveraging this approach, LongRoPE is able to expand the LLM context window to surpass 2 million tokens.

\subsection{Problem Formulation}

The presence of these two non-uniformities significantly expands the solution space and introduces additional challenges to optimization efforts. To manage this complexity, we recast the optimization of multidimensional non-uniform position interpolation as a search problem.

Consider a LLM configured for a context window size of $L'$, handling input documents $\mathbf{X}$ such that each element $\mathbf{x} \in \mathbf{X}$ exceeds $L'$ tokens in length. We represent the original rotary angle for the $i^{th}$ dimension within the RoPE embedding at token position $n$ as $\frac{n}{\beta^i}$. Based on this, we formally define the optimization problem as:

\begin{equation}
\begin{gathered}
		\label{eq:problem}

\underset {\mathbf{x}\in\mathbf{X}; \, |\mathbf{x}|\ge L'}{ \operatorname {arg\,min} } \,  \mathcal{L}\, \left( \text{LLM} (\text{RoPE}, \mathbf{X})\right), \quad \text{where} 
\\
\scalebox{0.93}{
$\underset {\begin{subarray}{l} i=0,\cdots,\frac{d}{2}-1;\\ n\in[ 0,|\mathbf{x}|);\end{subarray}}{\text{RoPE}_(n)}= {\left[\cdots,\cos\left(\mathbb{I}(\hat{\lambda}_{i}, \hat{n})\cdot\frac{n}{\beta^i}\right), \sin\left(\mathbb{I}(\hat{\lambda}_{i}, \hat{n})\cdot\frac{n}{\beta^i}\right),\cdots\right] }$}\\
\text{with} \quad \mathbb{I}(\hat{\lambda}_{i},\hat{n})=\begin{cases}
1 &\text{if $n<\hat{n}$}\\
{\frac{1}{\lambda}_{i}} &\text{if $n\geq\hat{n}$}
\end{cases}
	\end{gathered}
\end{equation}

Here, we introduce a group of rescaling factors, $\mathbb{I}(\hat{\lambda}_{i},\hat{n})$, to address both varieties of non-uniformities. The terms $\hat{\lambda_i}$ and $\hat{n}$ represent, respectively, non-uniformity along RoPE dimensions and at specific token positions. More precisely, $\mathbb{I}(\hat{\lambda}_{i},\hat{n})$ is employed to adjust the rotational angle in the $i^{th}$ RoPE dimension; in this context, $\hat\lambda_{i}$ serves as the dimension-specific rescale factor, and $\hat{n}$ acts as the positional threshold for tokens. For token positions less than $\hat{n}$ (specifically, the first $\hat{n}-1$ tokens), the rescale factor $\hat\lambda_{i}$ is inactive and the standard RoPE rotary angle $\frac{n}{\beta^i}$ remains in effect. Once token positions reach or surpass $n \ge \hat{n}$, however, the scaling factor is subsequently incorporated.

For a given target context window size $L'$, our goal is to determine the set of optimal rescale factors ($\mathbb{I}(\hat{\lambda}_{0},\hat{n})$, $\mathbb{I}(\hat{\lambda}_{1},\hat{n})$, ...$\mathbb{I}(\hat{\lambda}_{i},\hat{n})$...) across the $1^{st}$ through $d^{th}$ dimensions of RoPE. The aim is that, with these adjusted RoPE parameters, the target $\text{LLM}$ will yield the lowest possible next-token prediction loss $\mathcal{L}$ (or perplexity) on input data $\mathbf{X}$ consisting of $L'$ tokens.

\subsection{Searching the Non-uniform Position Interpolation}

To address the issue presented in Eq.~\ref{eq:problem}, we propose a straightforward yet powerful approach that identifies optimal RoPE rescale factors, thereby leveraging the complex, multidimensional irregularities found within position embeddings.

\textbf{Search space}. Our method encompasses a broad search space that captures both types of non-uniformity. Table~\ref{tbl:searchspace} outlines the specifics of this space. More precisely, we permit the optimization to determine a distinct rescale factor for every dimension within the RoPE embedding. To streamline the construction of the search space, the optimization is performed over $\lambda_i$ and $\hat{n}$, rather than directly over $\mathbb{I}(\hat{\lambda}_{i},\hat{n})$; here, $\hat{\lambda}_{i}=1/\lambda_i$. As depicted in Table~\ref{tbl:searchspace}, the parameter $\lambda_i$ can take values starting from 1.0, representing pure extrapolation, up to $s\times1.25$, denoting an interpolation broader than that in PI, with increments of 0.01. The parameter $s$ stands for the desired expansion ratio of the context window.

$\hat{n}$ determines how many of the first token positions are preserved using their unaltered RoPE embeddings, foregoing position interpolation. In practice, we select $\hat{n}$ from the set \{0, 1, 2, 4, 8, 12, 16, 20, 24, 28, 32, 64, 128, 256\} during the search process. If $\hat{n}=0$, then every token position is assigned the optimized rescale factors.

\textbf{Evolution-based search}. The search space outlined in Table~\ref{tbl:searchspace} encompasses a vast array of positional interpolation options, making exhaustive exploration highly impractical. For instance, utilizing a $s=4\times$ extension results in $400^{128/2}\times14=4\times10^{167}$ possible configurations. As the extension ratio increases, the number of combinations in the search space grows exponentially. To navigate this complexity, we adopt an evolution search strategy~\cite{guo2020single} and incorporate two optimization approaches that significantly enhance the efficiency of the search process. The complete search method is described in Algorithm~\ref{alg:search}.

\textit{Optimized initial population generation}. Rather than relying solely on random initialization for a population of $P$ rescale factors, we explicitly include the three RoPE rescale factors associated with PI, NTK, and YaRN as part of the initial pool. The other $P$-3 individuals are then created by introducing random mutations to these three rescale factors, each mutation occurring with probability $p$.

\textit{Monotonically non-decreasing constraint}. Once the initial population is produced, we evaluate the LLM perplexity associated with each candidate. To do this, the specific RoPE rescale factors are applied to the target LLM, and perplexity is measured on the input $\mathbf{X}$. The best-performing $k$ individuals are selected as parents for subsequent generations. Nevertheless, due to the enormous search space, straightforward mutation and crossover steps may frequently yield suboptimal candidates, resulting in superfluous perplexity evaluations. This inefficiency is exacerbated as $L'$ grows, since each perplexity computation incurs a considerable computational cost.

To this end, we enforce a constraint of non-decreasing monotonicity on the set of sampled RoPE rescaling factors such that $\lambda_i\le \lambda_{i+1}$ holds. RoPE variants that violate this criterion are excluded from use during LLM perplexity testing, thereby substantially narrowing the parameter search space. In detail, we stipulate that $\lambda_i$ must increase monotonically with respect to the RoPE dimension index $i$ (where $i=0,\ldots,63$). This monotonic relationship is supported by NTK theory~\cite{jacot2018neural,tancik2020fourier,ntk}, which implies that dimensions corresponding to higher frequencies (i.e., lower $i$) need minimal interpolation (lower $\lambda_i$), whereas dimensions with lower frequencies (higher $i$) can tolerate greater degrees of interpolation (higher $\lambda_i$).

\begin{algorithm}[t]
	\small
	\caption{The search algorithm}
	\textbf{Input:} target LLM, input samples $\mathbf{X}$, population size $P$, mutation size $N_1$, crossover size $N_2$, max iterations $\mathcal{T}$, mutate probability $p$\\

	\begin{algorithmic}[1]
		\label{alg:search}
		\STATE Initialize Top-k as empty;	
		\STATE Set $\text{P}_0$ = \textit{Initialize\_population\_with\_optimization} ($P$, $\mathbf{X}$, $p$); 
		\FOR{$i$ from $1$ to $\mathcal{T}$}
		\STATE \textit{Compute\_perplexity}(LLM, $\text{P}_{i-1}$, $\mathbf{X}$);
		\STATE  Top-k $\gets$ \textit{Update\_Topk}(Top-k, $\text{P}_{i-1}$);
		\STATE$\text{P}_{mutation}$ = \textit{Mutation\_with\_mono\_constraint}(Top-k, $N_1$, $p$); 
		\STATE $\text{P}_{crossover}$ = \textit{Crossover\_with\_mono\_constraint}(Top-k, $N_2$);
		\STATE $\text{P}_i$ = $\text{P}_{mutation}$ $\cup$ $\text{P}_{crossover}$ $\cup$ Top-k;
		\ENDFOR
		\STATE Output the individual with the minimum perplexity value in Top-k;
	\end{algorithmic}
\end{algorithm}

\textbf{8$\times$ extension without fine-tuning}. Through evolutionary search, our approach successfully discovers non-uniform RoPE rescaling factors, selectively maintaining crucial dimensions and positions in order to reduce information loss caused by interpolation. As shown in Fig.\ref{fig:non-ft}, this strategy enables LLaMA2's context window to be expanded from 4k to 32k tokens, all without requiring fine-tuning. Conversely, prior techniques like PI as well as non-uniform NTK and YaRN lead to a sharp increase in perplexity beyond a 2$\times$ context extension.

\subsection{Extending LLM Context Window to 2048K}
\label{sec:extendto2048k}

\textbf{Progressive extension to 2048k}. In this section, we present our approach for expanding the context window of pre-trained LLMs from the standard 4k to beyond 2048k. Our experiments show that by utilizing non-uniform positional interpolation, it is possible to attain up to an 8$\times$ increase without the need for fine-tuning. For more substantial extensions (such as 512$\times$), fine-tuning becomes essential. One strategy involves identifying appropriate RoPE rescaling factors suitable for the target 2048k window, followed by subsequent fine-tuning. However, this approach is hindered by the immense computational resources required. Additionally, our findings indicate that fine-tuning LLMs at such high extension ratios presents significant difficulties (refer to Appendix).

Fortunately, LongRoPE demonstrates efficacy for both base models and those that have undergone extended fine-tuning. Accordingly, we propose a progressive and efficient strategy that attains the desired 2048k target through only 1k fine-tuning iterations, all conducted within a training sequence length of 256k.

$\Diamond$ \textit{Exploring 256k context window expansion for pre-trained LLMs via LongRoPE search}. Using LLaMA2 as a case study, we perform a search with target context window sizes of 128k and 256k tokens. At this stage, the extension ratios reach 32$\times$ and 64$\times$, accordingly.

$\Diamond$ \textit{Fine-tuning to 256k}. Next, we adapt the pre-trained LLM to support a context window of 256k through further fine-tuning. Initially, we perform 400 fine-tuning steps on LLaMA2 with RoPE rescaled factors set for 128k. Following this, we substitute the RoPE rescaled factors with those corresponding to 256k at the resulting checkpoint and continue fine-tuning for an extra 600 steps. This staged approach demonstrates superior efficiency compared to directly fine-tuning the model to accommodate 256k.

$\Diamond$ \textit{Scaling fine-tuned extended LLM to 2048k using LongRoPE search}. In the final step, we apply an additional search procedure to the previously fine-tuned LLM with a 256k context length. This approach enables us to achieve a remarkably large context window of 2048k, eliminating the need for any subsequent fine-tuning. The total increase in length represents an extension factor of $512\times$.

\textbf{Recovery in shorter context windows}. When expanding the context window to a very large size such as 2048k, a decline in performance is observed within the initial context window. This degradation is a recognized limitation of positional interpolation~\cite{pi}: embedding positions from the original context window into a compressed segment of the higher-dimensional space leads to reduced language model effectiveness. Specifically, applying a 512$\times$ expansion means that the original 4k context window's positional encodings are densely packed, which exacerbates this issue.

To address this issue, we conduct an additional round of evolution search on the expanded LLM, tuning the RoPE rescale factors specifically for shorter context windows (such as 4k and 8k tokens). For these shorter lengths, we limit the upper bound for the searched $\lambda$ values, as they require less positional interpolation. At inference time, the LLM automatically applies the appropriate RoPE rescale factors as needed.

 \section{Experiments}

\label{sec:exp}
\subsection{Setup}
\textbf{Evaluation Tasks and models}. Our experiments involve integrating LongRoPE into LLaMA2-7B and Mistral-7B, followed by assessments across three dimensions: (1) measuring the perplexity scores of these extended-context LLMs when processing lengthy documents; (2) evaluating their capabilities in the Passkey retrieval task, which tests how effectively a model can identify a straightforward passkey hidden among large amounts of distracting text; and (3) assessing model performance on established LLM benchmarks confined to a context window of 4096 tokens.

\textbf{Fine-tuning}. For LLaMA2, fine-tuning is performed with a global batch size of 32 and an initial learning rate set at 2e-5, applying a linear decay schedule. The model undergoes 400 training iterations on the Redpajama~\cite{redpajama} dataset, which is partitioned into 128k-token segments and enclosed with BOS and EOS tokens. After completing these steps, an additional 600 training iterations are run from the same checkpoint to extend the context window to 256k tokens. For context length of 128k, training is distributed across 8 A100 GPUs using the system described in~\cite{cube}; 256k context window experiments require scaling up to 16 A100 GPUs. Regarding Mistral, training utilizes a fixed learning rate of 1e-6 and a global batch size of 64. For both the 128k and 256k variants, we adopt the procedure detailed in YaRN~\cite{yarn}, conducting 400 training steps on the Together Computer's Long-Data Collections~\cite{mistraldata} using a 16k sequence length, with computation performed on 4 A100 GPUs.

\textbf{Search}. When targeting window sizes up to 256k, our settings are as follows: $P$=64, $N_1$=$N_2$=16, $p$=0.3, $\mathcal{T}$=40, selecting the top 32 candidates for mutation or crossover at each iteration. Perplexity assessment is based on five randomly selected validation samples from PG19, ensuring each meets the minimum length corresponding to the target context. For window sizes exceeding 512k, the values for population, mutation, and crossover are reduced by half. In this case, perplexity is evaluated on three randomly chosen validation samples from the Pile-Books3~\cite{pile} dataset.

\textbf{Baselines}. To achieve a 2048k context window, we conducted fine-tuning on models pre-trained with context windows of 128k and 256k. The resulting models are referred to as LongRoPE-2048k (ft=128k) for LLaMA2 and LongRoPE-2048k (ft=256k) for Mistral. Our evaluation includes a comparison between these four models and leading context window extension methods, specifically open-source large language models that have undergone fine-tuning post-positional interpolation via PI, NTK, and YaRN approaches. Among these baselines are Together-32k~\cite{together}, Code LLaMA~\cite{codellama}, LongLoRA-full-FT-100k~\cite{longlora}, as well as YaRN-LLaMA and YaRN-Mistral~\cite{yarn}.

\subsection{Main Results}

\label{sec:mainresult}
\textbf{Long sequence language modeling within 256k}. Our initial analysis involves benchmarking against leading extended LLMs for sequence lengths up to 256k tokens. To showcase the robustness of our approach, we employ two separate datasets: the test splits from Proof-pile~\cite{pg19} and PG19~\cite{pile}. We report perplexity scores across multiple context lengths using a sliding window of 256 tokens. The full set of 100 test documents from PG19 is utilized in the evaluation. For Proof-pile, consistent with the methodology of YaRN~\cite{yarn}, we randomly sample 10 test cases, each having a minimum sequence length of 128k tokens.

Table~\ref{tbl:proof-pile} and Table~\ref{tbl:pg19} present a comparison of the perplexity values obtained by LLaMA2 and Mistral when augmented using various interpolation strategies on the Proof-pile and PG19 datasets, respectively. Two primary findings emerge from these results: \textbf{(1)} our extended models exhibit a consistently declining perplexity pattern as the evaluation length increases from 4k up to 256k, demonstrating their enhanced capacity to utilize longer contexts. \textbf{(2)} Despite operating with a context window that is 16$\times$ larger—a scenario that typically poses challenges in retaining performance for shorter contexts—our LongRoPE-2048k models surpass the latest baselines when assessed at a context length of up to 256k.

\textbf{Language modeling for sequences exceeding 2000k tokens}. To assess performance on very long document sequences, we utilize the Books3~\cite{pile} dataset. For the sake of evaluation efficiency, 20 books are randomly chosen, each containing more than 2048k tokens. A sliding window approach of 256k is then applied during evaluation.

Table~\ref{tbl:books3} illustrates that LongRoPE enables the extension of the context window for both LLaMA2-7B and Mistral-7B to 2048k tokens, while maintaining perplexity on par with or better than competing methods at shorter window sizes between 8k and 128k. We further note substantial disparities in performance between the 2048k variants of LLaMA2 and Mistral. Specifically, Mistral outperforms baseline models at reduced context lengths, yet its perplexity surpasses 7 once the context exceeds 256k. LLaMA2, conversely, behaves as anticipated: its perplexity improves as context window size increases, aside from minor rises at the 1024k and 2048k mark. Notably, with LLaMA2, the LongRoPE-2048k variant achieves superior results when fine-tuned with a window length of 256k compared to 128k, which is attributed to utilizing a smaller secondary extension ratio (8$\times$ vs. 16$\times$). On the other hand, Mistral attains better outcomes when fine-tuned at a window size of 128k. This discrepancy arises chiefly because, for Mistral’s 128k and 256k fine-tuning, we apply the YaRN protocol, limiting the training sequence to 16k, which subsequently constrains Mistral’s capability to extend its context window in the post-fine-tuning stage.

\textbf{Passkey retrieval}. Next, we analyze how well the model can leverage its context window for generation tasks. We adopt the synthetic passkey retrieval benchmark introduced by ~\cite{passkey}, where the model must extract a concealed random five-digit number (the passkey) from within an extended document. Specifics regarding the prompt template are available in the appendix. For evaluation, we conduct 10 runs, each time embedding the passkey at a location randomly selected with uniform probability over the evaluation context's length.

Fig.~\ref{fig:passkey} presents a comparative analysis of retrieval accuracy against baseline models. While previous approaches experience a steep decline in accuracy, falling to zero beyond 128k tokens, our LongRoPE-LLaMA2-2048k (ft=256k) achieves robust performance on the difficult task of extracting a passkey across up to two million tokens, consistently securing high retrieval accuracy ($\ge$90\%) from 4k through 2048k. Similarly, LongRoPE-Mistral-2048k (ft=128k) preserves perfect (100\%) accuracy up to 1800k tokens before declining to 60\% at 2048k, which corresponds with the observations in Table~\ref{tbl:books3}, where a minor increase in perplexity occurs at the 2048k context length.

\textbf{Evaluation on established benchmarks within the native context window}. We assess LongRoPE-2048k models on the original context window by utilizing the Hugging Face Open LLM Leaderboard~\cite{huggingfaceboard}, conducting both zero-shot and few-shot experiments. Specifically, we apply 25-shot settings for ARC-Challenge~\cite{allenai:arc}, 10-shot configurations for HellaSwag~\cite{zellers2019hellaswag}, 5-shot scenarios for MMLU~\cite{mmlu}, and employ a 0-shot format for TruthfulQA~\cite{lin2021truthfulqa}.

Table~\ref{tbl:commonsense} demonstrates that our models perform on par with existing systems on the initial benchmark—intended for limited context windows—and exceed Mistral's results on TruthfulQA by an improvement of +0.5\%. Although LongRoPE-LLaMA2-2048k, which underwent fine-tuning at 256k, experiences a marginally greater decline in performance, its scores for the majority of tasks stay within acceptable bounds.

\subsection{Ablation Results}

\textbf{Effectiveness of the second positional interpolation}. Within our progressive extension framework, a second stage of non-uniform positional interpolation is applied to the fine-tuned, extended LLMs, utilizing our search algorithm. To assess its utility, we experimented with our fine-tuned LLaMA2-256k model, which we further extended to 512k, 1024k, and 2048k sequences using both PI and YaRN methodologies. As demonstrated in Table~\ref{tbl:secondsearch}, the non-uniform positional interpolation method we employ maintains stable perplexity values. By contrast, perplexity in models extended with PI and YaRN escalates markedly as the extension ratio increases.

\textbf{Effectiveness of recovery at shorter context lengths.} In order to address performance degradation at reduced context lengths, we re-optimized the RoPE factors for LongRoPE-2048k utilizing our search algorithm. In this process, we specifically restricted the maximum permissible scale factors during the search phase to promote reduced interpolation for shorter contexts of 4k and 8k tokens. Table~\ref{tbl:shortperformance} presents the perplexity outcomes for LongRoPE-LLaMA2-2048k on the Proof-pile dataset at 4k and 8k lengths, as well as the mean LLM benchmark accuracy. These findings indicate a pronounced improvement in performance for shorter context lengths.

\textbf{Analysis on the two forms of non-uniformities}. To disentangle the effects of the two types of non-uniformities, we conduct an ablation study to assess their individual contributions to model performance. Two experimental scenarios are explored: (i) for LLaMA2-7B, we extend the context length to 16k and 32k tokens, comparing methods such as PI, RoPE dimension-only optimization, and optimization involving both types of non-uniformities; (ii) for our fine-tuned LLaMA2 model at 256k length, we further extend it to 2048k using analogous approaches. Evaluation of perplexity is performed without any additional fine-tuning. As indicated in Table~\ref{tbl:searchdimension}, optimizing the RoPE dimension non-uniformity yields a marked decrease in perplexity relative to PI's linear interpolation method. Adjusting non-uniformity in the token position demonstrably enhances model accuracy at 16k and 32k tokens, but such improvements are not as evident at 2048k, likely attributable to the extraordinary sequence length. Moreover, merely retaining initial tokens without applying interpolation becomes ineffective, and we note this as a potential avenue for future study.

\section{Related Works}

In addition to methods based on position interpolation, this section discusses related works of other approaches. 

\textbf{Retrieval-based} methods rely on an external memory component to store extended historical context and employ retrieval mechanisms to fetch pertinent documents during inference~\cite{longllama,wang2023augmenting,borgeaud2022improving}. Such methods generally require direct and explicit adjustments to the underlying LLM architectures. By contrast, our approach introduces only minimal changes to positional embeddings, resulting in a more streamlined and efficient solution. Furthermore, our method accommodates a broader range of long context applications, including long-form document summarization and few-shot learning, beyond retrieval-specific scenarios.

\textbf{Attention-based context window extensions}. In addition to extending input context through positional embedding interpolation, several studies have increased the context window of LLMs within their existing architecture by adapting attention mechanisms~\cite{han2023lminfinite, streamingllm,parallelwindow}. Central to these methods is the use of innovative attention masks to address the issue of attention explosion from newly introduced positions. These strategies function synergistically with approaches based on positional interpolation.

\textbf{Fine-tuning based approaches} investigate strategies for adapting pre-trained LLMs by adjusting their position embeddings to handle longer sequences. For instance, Code LLaMA~\cite{codellama}, LLaMA2 Long~\cite{xiong2023effective}, and ScaledRoPE~\cite{liu2023scaling} employ an enlarged base value for RoPE, subsequently fine-tuning the models on the desired target context length. Unlike these, our method is adaptable across multiple target lengths and is capable of supporting sequences exceeding 2M tokens. More recently, since fine-tuning at extended context lengths (greater than 128k) incurs heavy GPU requirements, solutions such as LongLoRA~\cite{longlora} and PoSE~\cite{zhu2023pose} have been designed to reduce this computational burden. It should be noted that our approach is complementary to these efficient fine-tuning techniques.

\section{Conclusion}

In this study, we introduce LongRoPE, a technique that significantly increases the context length of LLMs to an extraordinary 2048k tokens, all while preserving model performance within their original, shorter context range. Our approach leverages two types of non-uniformities in RoPE positional encodings, identified through an efficient evolutionary search process. This method yields dual advantages: it supplies optimal initialization for subsequent fine-tuning and supports an 8$\times$ extension of the context window without the need for any fine-tuning. On this foundation, we implement a progressive extension scheme by fine-tuning LLMs on sequences of length 256k, eventually achieving a 2048k context window, again without additional fine-tuning steps. Comprehensive experiments demonstrate the efficacy of LongRoPE. We anticipate that our LongRoPE-2048k models will unlock numerous long-context tasks and serve as a catalyst for future investigation in this area.

\section*{Broader Impacts} This research aims to further developments within the Machine Learning discipline. While our work could have various implications for society, we do not identify any particular issues that warrant explicit emphasis in this context.

	\balance

\end{document}